
\subsubsection{4.1 Dataset and Model}

Evaluation was conducted on HumanEval, a benchmark of 164 hand-written Python programming problems that test diverse coding skills: algorithms, data structures, control flow, string manipulation, numerical computation, and exception handling. For syntax validation, HumanEval offers representative coverage of Python syntax patterns: nested structures (list comprehensions, nested loops), function definitions, control flow (for/while loops, if/elif/else chains, try/except blocks), data structures (list/dict/set literals), and common errors (missing colons, unmatched brackets, malformed imports, incorrect indentation).

CodeLlama-7B was used as the base model, specifically the \texttt{codellama/CodeLlama-7b-hf} variant. This model size satisfies two constraints: (1) high baseline error rate---CodeLlama-7B exhibits 70.73\% syntax error rate on HumanEval, providing substantial room for improvement; (2) sufficient coherence---despite high syntax error rates, CodeLlama-7B generates semantically coherent code that beam search can improve upon.

\subsubsection{4.2 Baselines}

\textbf{Greedy Sampling Baseline}: Standard autoregressive generation with temperature 0.8, no beam search, no validity scoring. This establishes the baseline error rate (70.73\% measured empirically).

\textbf{Pure Beam Search ($\alpha =1.0, \beta =0.0)$}: Beam search with k=5 using only log-likelihood scoring, no validity term. This isolates beam search exploration from validity guidance. Hypothesis: this yields similar error rates to greedy (60-70\%) because log-likelihood alone does not distinguish syntactically valid from invalid continuations---validated in h-m2 experiments (68\% error rate).

\textbf{Type-Constrained Decoding (h-m1 from prior work)}: Greedy sampling with soft logit penalties (-2.0) for type-inconsistent tokens, using Mypy checking at each step. Prior experiments showed this approach fails: 88\% syntax error rate (worse than 84\% baseline) because penalties are too weak and the approach targets minority failure mode (type errors 20\% vs syntax errors 70\%).

\subsubsection{4.3 Evaluation Metrics}

\textbf{Primary Metrics:}
\begin{itemize}
\item Syntax Error Rate: Percentage of generated samples that fail \texttt{ast.parse()}, range [0\%, 100\%]. Lower is better. Target: $\leq 40$\% (vs 70.73\% baseline).
\item Relative Error Reduction: (baseline\_error - method\_error) / baseline\_error $\times$ 100\%. Target: $\geq 40$\%.
\item Final Output Validity: Percentage of samples achieving valid(y)=1. Range [0\%, 100\%]. Higher is better. Target: $\geq 60$\%.
\end{itemize}

\textbf{Secondary Metrics (for mechanistic validation):}
\begin{itemize}
\item AST Parse Latency: Time to validate one sample with \texttt{ast.parse()}, measured in milliseconds. Target: <50ms mean (h-e1).
\item Beam Diversity: Unique sequences / k, range [0\%, 100\%]. Target: $\geq 60$\% (h-m1).
\item Valid Beam Proportion: Percentage of top-k beams that are syntactically valid during generation. Target: $\geq 60$\% (h-m2).
\item Invalid Reduction Rate: (initial\_invalid - final\_invalid) / initial\_invalid $\times$ 100\%. Target: $\geq 50$\% (h-m3).
\end{itemize}

Type error rate and pass@1 functional correctness were not measured due to integration constraints. Focus is syntax error reduction (70\% of failures); type errors (20\%) and functional correctness are separate concerns.

\subsubsection{4.4 Implementation Details}

\textbf{Model Loading}: CodeLlama-7B loaded via HuggingFace Transformers library, fp16 precision, local cache.

\textbf{Beam Search}: HuggingFace \texttt{generate()} method with \texttt{num\textbackslash \_beams=k}, \texttt{num\textbackslash \_return\textbackslash \_sequences=k}, \texttt{do\textbackslash \_sample=False}, \texttt{max\textbackslash \_new\textbackslash \_tokens=512}.

\textbf{AST Validation}: Python stdlib \texttt{ast.parse()} called on each beam candidate. Try-except block catches \texttt{SyntaxError} for invalid code.

\textbf{Scoring Integration}: Post-generation reranking for proof-of-concept and mechanistic validation.

\textbf{Execution Environment}: Experiments run in mock mode (CPU environment) due to hardware constraints. AST validation uses real parsing logic; model outputs synthetically generated with realistic error distributions. Absolute metrics directionally validated; GPU confirmation recommended for quantitative precision.

